\section{One finite construction for every real time and height}
\label{sec:all-time}
The effective positive table selects a finite $T_+$. We must next pay the
whole interval between the earlier switch construction and that overlap.
The constants below explicitly depend on the now fixed $b=T_++1$.
They do not enter the already chosen $C_+$ or $T_+$.

\subsection{A compact exterior table with an arbitrary fixed upper time}
For $J_b=[63,b]$, $b\ge128$, put
\[
 \Lambda_b=2^{20}(1+b)^4,\quad C_0(b)=12b^{3/2}+7,\quad
 E_b^{\rm out}=2^{60}(1+b)^4,\quad D_b^{\rm out}=2^{32}(1+b)^4.
\]
Repeating the complete contour proof on $x\ge e^{\Lambda_b}$ gives
$x/2\le|p|\le x$, $t\log x/5\le\Re(p-w)\le t\log x/3$,
$|\Im(p-w)|<t/2$ and $|B-1|\le b/x$.
The central error is at most $12b^{3/2}/x^2$. The remaining finite and
infinite contour pieces are each at most $x^{-3}$: their common exponent
is bounded by
\[
 -63\log^2x/100+(2b/3)\log x+6\le-\tfrac12\log^2x,
\]
and the complete arithmetic tail is at most
$48e^{b/2-63\log x/10}$. These inequalities hold at $\log x\ge\Lambda_b$.
Thus on strip five $|H-P|\le C_0(b)A_\Sigma/x^2$.
Complete Cauchy circles give, on strip $9/2$ and
$x\ge2e^{\Lambda_b}$,
\begin{equation}
 |\partial_z^j(H-P)|\le
 8C_0(b)j!4^j\log^j(2+x)A_\Sigma/x^2\quad(1\le j\le3).
 \label{eq:at-compact-jets}
\end{equation}
The exact saddle derivatives also give the earlier phase and collar
constants, $|X'(t)|\le8$, and residual bounds
$|E_P|\le D_b^{\rm out}A_\Sigma/x^3$,
$|(E_P)_z|\le D_b^{\rm out}\log(2+x)A_\Sigma/x^3$.
The velocity estimate uses the real-part cancellation in
\eqref{eq:ua-speed}, with $t/x^2\le b/x^2$.

The full complex log-gamma expression supplies the sharper carrier
estimate needed later:
\begin{equation}
 \left|\log|S_t((1\pm iz)/2)|+\pi x/8\right|
 \le A_b^{\rm amp}(1+\log x)^2,\qquad
 A_b^{\rm amp}=2^{16}(1+b)^2.
 \label{eq:at-carrier}
\end{equation}
Indeed $\Re p\le3+b\log x/2$, the imaginary shift is at most $b/2$,
the angular correction is bounded by a constant times $\Re p+b$,
and the real saddle exponent by $b\log^2x/9+b/4$.
The logarithmic, polynomial, determinant and sector-remainder terms all
fit the displayed coefficient. At this onset it is smaller than $x/8$.
Consequently $e^{-x}\le|S_t|$ on the full strip and
$|S_t((1\pm ix)/2)|\le e^{-x/4}$ on the real axis.
This is a new fixed-$b$ table, not an extrapolation of the time-$97$ table.

\subsection{Regulator, moments and means on the entire interval}
The real source saddle satisfies
\[
 1/4<L_s<b+1,\quad36<\kappa<K_b:=8b(b+1)+36,
 \quad0<\kappa'<K_b':=16(1+b)^2.
\]
These follow by bracketing its equation at $1/4,b+1$ and using its
exact derivative. With $W_b=16(1+b)^3$, the source-saddle formula gives
$e^{-3}K_b^{-1/2}\le M_0\le e^{W_b}$.
Write $r=a-\beta x^2$ as in \eqref{eq:ua-switch-r}, and set
\begin{align*}
 L_{+,b}&=W_b+3+2\log(2+K_b),&
 L_{0,b}&=|\log\eta|+4+\tfrac14\log(1+b^2),\\
 A_b&=4[1+L_{+,b}+L_{0,b}+(1+b)^2+K_b'],&
 B_b'&=1/36+K_b'/36^2,\\
 Q_b&=16\eta^2+e^{2W_b},&
 F_b&=\max(4,2K_bB_b').
\end{align*}
Then $|a|,|a'|\le A_b$, $e^a\le Q_b$, $\beta\le1/36$,
$|\beta'|\le B_b'$, and
\begin{equation}
 |\beta'|\le F_b\beta[1+\log^2(1/\beta)],\qquad
 E_b=\min\left(\frac{4\eta e^{-4}}{(1+b^2)^{1/4}},
                    e^{-3}K_b^{-3/2}\right),\quad
 \sqrt q\ge E_be^{-R^2/72}\quad(|x|\le R).
 \label{eq:at-regulator}
\end{equation}
At $\beta=0$ its derivative is zero. On $[64,97]$ use the proved switch
flatness; on $[97,b]$, $\chi\ge1/2$ gives $\beta\ge1/(2K_b)$,
which proves the enlarged bound. The term
$\chi'(\log\epsilon_+-\log\epsilon_b)$ is retained in $a'$.

For full supports in $|x|\le R$, a sufficient physical coefficient
table is
\[
 A_t=A_b+B_b'R^2,\quad R_1=R/18,\quad R_2=1/18,\quad
 T_1=2B_b'R,\quad T_2=2B_b',\quad T_h=4B_b'R^2,
\]
with $c_r=1/18+R^2/324$ and $b_r=A_t+c_r/2$.
On exterior supports the same two cases as the switch proof apply.
If $\beta u\le16$, flatness gives $|r_t|=O_b(u\log^2(2+u))$.
If $\beta u>16$, the relative physical regulator satisfies
$\Lambda\le\sqrt{Q_b}e^{-2u}$ at the final onset. These cases split
positive allowances only; no indicator is differentiated.

For complete moments let $V_b=b+16$ and
\begin{equation}
 B_j(b)=40\left[\frac{V_b^{j+1}e^{bV_b^2+13V_b}}{j+1}
                     +\frac{\Gamma((j+1)/2)}2\right],\quad 0\le j\le3.
 \label{eq:at-moments}
\end{equation}
The source bound $\Phi(u)\le40e^{9u-\pi e^{4u}}$ and
$bu^2+13u+u^2\le u^3\le\pi e^{4u}$ for $u\ge V_b$ prove this
by splitting the integral at $V_b$. Thus these moments cover every
$t\le b$, all real heights, and strip four.

To specify the canonical mean coefficient without an unknown compact
supremum, put
\begin{align*}
 S_b&=1+Q_b+A_b+F_b+B_b'+E_b^{\rm out}+D_b^{\rm out}+\eta^{-2},\\
 R_b^*&=\left\lceil e^{2\Lambda_b+2S_b}\right\rceil,
 &e_b^*&=E_be^{-(R_b^*+8)^2/72},\\
 C_{H,b}&=\max\left(2^{100},
 2^{30}\left[\frac{B_0(b)+B_2(b)}{2e_b^*}+(R_b^*+8)/18+1\right]\right).
\end{align*}
Inside this fixed rectangle use the positive floor and moments; outside
use the paid units, integrable poles and weighted regulator estimates.
Moving remaining derivatives to $K$ gives
\begin{equation}
 |\partial^\gamma U_H|\le C_{H,b}[1+\log(2+|x|)],\quad1\le|\gamma|\le3.
 \label{eq:at-means}
\end{equation}
The actual joined mean has this cap plus one after its auxiliary
primitive rows are paid. This order avoids circular covariance bounds.
The compact $C_{H,b}$ is never used in the late-time $C_+$.

\subsection{One gap, one compact budget, one source}
Choose a phase level at time $64$ above a finite onset $U$:
$n=\lceil\Theta(64,U)/\pi\rceil$, $\Theta(64,X_0)=n\pi$.
Follow this same level throughout $J_b$. The velocity bound gives
$|X(t)-X_0|\le8\max(1,b-64)$.
Define
\[
 D_{\rm mov}=8(b+1)+32,\quad a_X=X_0-D_{\rm mov},\quad
 R_c=X_0+D_{\rm mov},\quad L=\log(2+R_c),\quad \sigma=(64L)^{-1}.
\]
All collar, source and averaging buffers fit inside $R_c$; the model
and actual gap floors are respectively $1/2$ and $1/4$.
Substitute
\[
 \mu_o=E_b^{\rm out}/a_X^2,\quad\mu_s=a_X^{-2},\quad
 d=d_1=D_b^{\rm out}/a_X^3
\]
in the spatial and moving ledgers, retaining every actual time derivative.
The same finite maps as in Section~\ref{sec:uniform-assembly} give the
following sufficient coefficient and nonsource envelope:
\begin{align}
 C_b={}&2^{8192}
 [1+C_{H,b}+Q_b+A_b+B_2(b)+1+F_b+B_b'
                 +E_b^{\rm out}+D_b^{\rm out}+D_{\rm mov}]^{64},\notag\\
 B_b(X_0)={}&C_b\{L^8/a_X+(1+R_c)^{20}e^{-a_X}\}.
 \label{eq:at-compact-budget}
\end{align}
The worst time-compatibility row is $O_b(\mu uL^3)$; all other power
tails decay at least as fast. The time length enters $D_{\rm mov}$,
not an accumulated count of charts. The finite maps use fewer than
sixteen powers of the displayed inputs; the sixty-fourth power and
$2^{8192}$ dominate their profile, kernel and Leibniz coefficients.

For the final positive band use $C_b^+=(1+c_b)C_b$ and the corresponding
$B_b^+$. Given a reserve $\delta$, a sufficient binary onset is $U=2^m$,
where
\begin{align}
 m\ge{}&16384+\lceil\log_2(4R_b^*+4D_{\rm mov}+1024)\rceil
                   +\lceil8\log_2(2+C_b^+/\delta)\rceil,\notag\\
 m\ge{}&16384+\lceil8\log_2(2+2^{64}K_{\rm br}/\delta)\rceil.
 \label{eq:at-onset}
\end{align}
The phase inverse enlarges $U$ by less than one at this onset.
Then $a_X\ge2^{m-1}$, $L\le m+3$,
$C_b^+\le\delta2^{m/8}$ and $(m+3)^8\le2^{m/2}$;
the exponential term is smaller still. Thus $B_b^+<\delta/4$, and
all branch, collar and scale gates hold. The second inequality pays
the earlier pre-switch bridge.

The original sufficient core source demand is explicit. Set
$e_c=E_be^{-R_c^2/72}$, $\rho=\min(\delta/4,B_b^+)$ and
\[
 \varepsilon_c=e_c^2\min_{D_r>0}\frac{\rho}{4D_r},
 \qquad
 e_{\min}=\frac{T_+^{3/2}}{[1+R_c+b\log(2+b+R_c)]^2}.
\]
Here $D_r$ are the finite compact-transfer row coefficients obtained
from \eqref{eq:at-regulator}--\eqref{eq:at-means}, including auxiliary
mean rows. With $T_-=\max(1,4K_-/\delta)$, choose
\begin{equation}
 \varepsilon\le\min\left\{1,\varepsilon_c,
 a_X^{-2}e^{-R_c-16},T_-^{-5/2},
 \frac{\delta}{4(1+K_{\rm mid}(T_-))},
 \frac{\delta}{2^{69}K_{\rm br}L^3},
 \frac{e_{\min}A_{\rm late}}{16}\right\}.
 \label{eq:at-source-minimum}
\end{equation}
All entries are finite positive numbers already specified. The clipped
constructor of Section~\ref{sec:effective-inputs} applies with upper time
$b$ after replacing its cutoff by
$U_s=\lceil b+4+\log(5+\log(32/\varepsilon))\rceil$ and
$F_s=U_s^4e^{bU_s^2+4U_s}$; its complete tail has slope at least
$\pi e^{4U_s}$ and exponent at most $-2e^{4U_s}$.
It gives one nonnegative source whose four jets have errors
$29\varepsilon/64$ for \emph{every} $t\le b$ and all real heights
on strip four. Section~\ref{sec:source-conditioning} improves this
sufficient tolerance and the source cutoff; the current prescription
already proves finite existence.

\subsection{The positive band and the global domain table}
On $[T_+,T_++1]$ use $v_+=b_*(t-T_+)$ and
$\lambda_+=v_++(1-v_+)\lambda$. If $B_c=(F_c)_{\bar z}$, then
\begin{align*}
 (\lambda_+)_t&=v_+'(1-\lambda)+(1-v_+)\lambda_t,&
 B_+&=(1-v_+)B_c,\\
 (B_+)_t&=(1-v_+)(B_c)_t-v_+'B_c,&
 ((B_+)_z)_t&=(1-v_+)((B_c)_z)_t-v_+'(B_c)_z.
\end{align*}
The heat terms are $\lambda_+E_P$,
$(1-v_+)[\lambda_x(P-G)_z+\lambda_{xx}(P-G)/4]$ and
$[(1-v_+)\lambda_t+v_+'(1-\lambda)](P-G)$.
On the band the additional compatibility variation in
\eqref{eq:tc-signed-heat} is $O(|v_+'|\mu L^2)$; area $O(1/L)$
makes its filtered cost $O(|v_+'|\mu L)$. This fits $B_b^+$.
The core-containing heat pieces are paired with the global weak test;
no holomorphic identity is applied to a sharply truncated core integral.
All old spatial commutators and the new compatibility variation remain.

\begin{theorem}[Finite all-real-time approximation]
\label{thm:all-time}
Fix $\eta>0$, $K=k_1k_1$ and $0<\theta\le1$. There are explicitly
specified finite parameters and one nonnegative finite source such that
the following smooth field controls every scalar of
\eqref{eq:ua-contract} within $\theta/2$, uniformly on
$\R_t\times\R_x\times\{|y|\le1\}$:
\begin{equation}
 F(t,z)=\begin{cases}
 N,&t\le-2T_-,\\
 (1-v_-)N+v_-G,&-2T_-\le t\le-T_-,\\
 G,&-T_-\le t\le63,\\
 G+v_0\lambda(P-G),&63\le t\le64,\\
 G+\lambda(P-G),&64\le t\le T_+,\\
 G+[v_++(1-v_+)\lambda](P-G),&T_+\le t\le T_++1,\\
 P,&t\ge T_++1.
 \end{cases}
 \label{eq:at-field}
\end{equation}
Here $v_-=b_*((t+2T_-)/T_-)$, $v_0=b_*(t-63)$ and
$N=C(-t)^{-1/2}e^{z^2/(4t)}$. The source and parameters may depend
on $\eta,\theta$. Equation residuals have the sums of their displayed
component allowances. Boundary traces inherit the stated derivative
orders and full buffers.
\end{theorem}
\begin{proof}
First choose $C_+,T_+$ by Section~\ref{sec:effective-positive}, then
$b=T_++1$ and $\delta=\theta/64$. This fixes the compact table.
Choose $T_-$ as above, then $m$, the phase level and gap, and finally
one source satisfying \eqref{eq:at-source-minimum}.
Section~\ref{sec:temporal-cover} pays the first four rows at reserve
$\delta$. The compact spatial proof pays the fifth by
$3B_b^+\le3\delta/4$. The sixth uses $\theta/8$ on its holomorphic
core and $B_b^+$ on its spatial bands; the last uses $\theta/8$.
Whole-support positive primitive allowances are added before the weak
identities and actual covariance mean product are evaluated. These
positive sums fit $\theta/2$. There is no multiplier for the number of
time rows, since only the field active at that time is compared.

Flat profiles give smooth endpoint matching. Before the late onset $P$
is used only on its paid exterior branch; its zero cutoff extends into
the inner region. The all-height late branch agrees with that branch by
uniqueness. The new compact table covers $97<t<T_++1$ explicitly.
No source coefficient uses an unknown zero of $H$ as data.
\end{proof}
At negative infinite time $F=N$, so the complete negative endpoint
errors vanish at their proved rates. At positive infinite time $F=P$:
static $U,W$ errors are at most $C_{\rm raw}e$ and the retained evolution
and covariance errors at most $C_+e\mathcal L^2\le C_+t^{-1/2}$,
uniformly in height. The model heat residual retains its separate $d$
scale. These are limits of approximation errors in physical units;
the density itself need not have a finite limit.

The construction is nonholomorphic on spatial bands and has nonzero heat
defects. Its algebraic density there is not asserted to count zeros.
It proves finite analytical existence for each positive target. It does
not yet give affordable source data, a working-precision theorem or a
continuous averaging certificate.
